\section{Methodology}
\label{sec:methodology}

\subsection{Study design and unit of analysis}
\label{subsec:design}

The planned study uses mixed methods: qualitative content analysis with a
frozen codebook, followed by non-parametric analysis of marker frequencies.
The target corpus consists of English-language user accounts concerning
Replika or Character.AI. The intended unit of analysis is the author, with
markers and word counts pooled across that author's posts before
normalization. The current synthetic test data do not contain author IDs;
therefore, treating one row as one author is a pipeline-test assumption and
must not be carried over as an empirical sampling fact.

\subsection{Measures and formalization}
\label{subsec:measures}

For author $i$, let $w_{ij}$ be the word count for post $j$ and let
$m_{ij}^{(k)}$ be the coded count for dependence marker $k$. Equation~
\ref{eq:dependence-rate} defines the planned emotional-dependence rate per
1,000 words:
\begin{equation}
  F_i = 1000\,
  \frac{\displaystyle\sum_{j=1}^{J_i}\sum_{k\in\mathcal{K}_F}m_{ij}^{(k)}}
       {\displaystyle\sum_{j=1}^{J_i}w_{ij}},
  \label{eq:dependence-rate}
\end{equation}
where $\mathcal{K}_F$ contains attachment, fear-of-loss, and daily-support
markers. For social-avoidance counts $a_{ij}$, the corresponding index is
defined in Equation~\ref{eq:isolation-index}:
\begin{equation}
  I_i = 1000\,
  \frac{\displaystyle\sum_{j=1}^{J_i}a_{ij}}
       {\displaystyle\sum_{j=1}^{J_i}w_{ij}}.
  \label{eq:isolation-index}
\end{equation}

Equations~\eqref{eq:dependence-rate}
and~\eqref{eq:isolation-index} define the
planned author-level rates. The marker
composition is specified in
\texttt{configs/study.yaml} and remains
subject to confirmation before analysis
of authentic data.

\subsection{Collection, coding, and analysis workflow}
\label{subsec:workflow}

The planned workflow is summarized in Fig.~\ref{fig:workflow}. The repository
configuration specifies a minimum post length of 100 words and other quality
filters. 
A second coder is planned for a random
10--20\% subsample, with Cohen's
$\kappa\geq 0.60$ as a preliminary
reliability target. Coding agreement will
be evaluated following established
content-analysis reliability principles~
\cite{krippendorffReliability2004}.
Disagreements will prompt codebook review.
For RQ2, the primary contrast is a one-sided Mann--Whitney
$U$ test between the less-than-one-month and more-than-three-month groups,
with Cliff's $\delta$ and a bootstrap 95\% confidence interval. Kruskal--
Wallis and Jonckheere--Terpstra tests are specified for secondary group and
trend analyses. For RQ3, the plan uses Spearman rank correlation and an OLS
model controlling for weekly session frequency, specified in
Equation~\ref{eq:rq3-model}:
\begin{equation}
  I_i = \beta_0 + \beta_1\,\mathrm{duration}_i
      + \beta_2\,\mathrm{sessions}_i + \varepsilon_i.
  \label{eq:rq3-model}
\end{equation}
The study is cross-sectional and observational, so results will be
interpreted as associations only.


\begin{figure}[t]
  \centering
  \begin{tikzpicture}[
    node distance=0.55cm,
    stage/.style={draw, rounded corners, align=center,
      text width=0.88\columnwidth, minimum height=0.62cm},
    flow/.style={-{Latex[length=2mm]}, thick}
  ]
    \node[stage] (corpus) {Public user accounts\\(planned authentic corpus)};
    \node[stage, below=of corpus] (filter) {Eligibility and quality filtering};
    \node[stage, below=of filter] (coding) {Codebook coding\\and reliability check};
    \node[stage, below=of coding] (rates) {Pool by author; normalize markers\\per 1,000 words};
    \node[stage, below=of rates] (tests) {RQ2 group tests and RQ3\\association analysis};
    \draw[flow] (corpus) -- (filter);
    \draw[flow] (filter) -- (coding);
    \draw[flow] (coding) -- (rates);
    \draw[flow] (rates) -- (tests);
  \end{tikzpicture}
  \caption{Planned content-analysis workflow. Synthetic repository data are
  used to validate code and do not instantiate the planned human corpus.}
  \label{fig:workflow}
\end{figure}

The project repository also contains synthetic benchmark experiments for
software performance. Those engineering measurements are separate from the
behavioral research questions and are not included in the analysis described
above.
